\chapter{Trajectory Output and Restart}\label{ch:output}

The product of a run is a set of trajectories: the recorded position, and any sampled properties, of every
particle through time. This chapter describes the two output backends and how to choose between them, the
variables written and their arrangement (output format~2, the format of release 1.0), the record time, the
provenance recorded with the trajectories, the checkpoint mechanism that allows a long run to be resumed and
continued in a chain of jobs, and the merging of a release split across processes. The authoritative
variable-by-variable description is \code{docs/io\_and\_time.md}, section~5.

\section{Two backends}\label{sec:out-backends}

The output format is chosen by the file extension of the output path. A \code{.zarr} path selects a chunked
Zarr store; any other selects a compressed NetCDF4 file (an unlimited record dimension, compression level 4).
The distinction is operationally significant for long runs. A NetCDF4 file is held under an exclusive write
lock for as long as it is open, so it cannot be inspected while the run proceeds; the key
\conf{output\_flush\_every} (command line \code{--flush-every}) mitigates this by buffering that many records
and then opening, appending, and closing the file, so that completed records become readable between flushes
and the file stays byte-identical. A Zarr store, by contrast, is lock-free and can be read while the run is
still writing. It buffers records in memory and writes one chunk of records at a time, so readers see it grow
in chunk-sized increments and a killed run loses at most one chunk of records; a checkpoint and the end of the
run flush the buffer. For multi-hour runs that one may wish to monitor, Zarr is the more convenient choice.

\paragraph{The Zarr layout.} Each variable is chunked in blocks of $n_t$ particles by $n_o$ records,
\begin{equation}\label{eq:out-chunks}
  n_t=\min(P,\ \num{500000}),\qquad
  n_o=\max\Bigl(1,\ \min\bigl(256,\ \bigl\lfloor 4\cdot10^{6}/(8\,n_t)\bigr\rfloor\bigr)\Bigr),
\end{equation}
so a chunk holds about \SI{4}{MB}. Up to half a million particles the particle axis is one block and a chunk
is one or more complete time columns; above that, a subset of particles reads only the blocks it touches.
Outputs written before this layout (August 2026) keep a single particle block.

\section{What is written}\label{sec:out-schema}

The output follows the CF conventions (the Climate and Forecast metadata conventions) for trajectories and
carries the global attribute \code{kelp\_output\_format} with the value~2; a file without it was written by an
earlier engine (\cref{sec:out-format1}). It is organized as a set of two-dimensional arrays indexed by
particle and by record, the \emph{trajectory} and \emph{time} dimensions (\cref{fig:output}), and both
dimensions have a coordinate variable of the same name. Every run writes the position --- longitude,
latitude and depth (positive down) --- the particle status as a coded flag (1 alive, 2 beached, 3 outside,
4 not released), the record time, a trajectory identifier, and, once per particle, the release position and
release time. The identifier is the particle's index in the whole release, so a shard of a sharded release
(\cref{sec:out-merge}) numbers its particles from its first global index. Several further variables are
written when they are present: the velocity as \code{vel\_u}, \code{vel\_v} and, in three dimensions only,
\code{vel\_w} (described below), the failure reason (\cref{sec:life-reason}), and one array per
sampled scalar field (\cref{ch:sampling}). Positions, release positions and horizontal velocities are always
written in geographic coordinates, un-rotated from any computational frame (\cref{ch:frame}). Where nothing
applies the arrays hold a fill value: a missing value (NaN) for positions, velocities, sampled fields and
release variables, and zero (no failure) for the failure reason; the rows of particles not yet released are
missing values. Positions carry the CF names and units (\code{longitude}, \code{latitude}, and
\code{depth} with \code{positive = "down"}), the status, failure reason and sampled fields name the
positions as their coordinates, and both backends write the same attributes. The one known departure from
CF is the order of the dimensions: CF wants the unlimited record dimension to be the outer one, while here
it is the inner one, so that the track of one particle is contiguous on disk.

A record is written after every step $g$ for which $g+1$ is a multiple of \conf{output\_freq}, counting steps
from the start of the run across any restarts. The release positions are therefore not part of the records ---
a particle first appears at the end of a step, at least one step after its release --- and are kept in the
separate release variables instead. The size of the output is proportional to
the number of particles times the number of records; for a large pulsed release this product, not the number
of live particles, sets the size, because the frozen particles occupy rows in the pre-allocated arrays. A run
with several realizations writes one file per realization, with the suffix \code{\_r}$N$.

\paragraph{The record time.} The time is stored in seconds since the first record of the forcing time axis
as it appears in the file, and the attributes of the time variable name that record's date as the epoch and
the forcing's calendar (for example \code{noleap} or \code{360\_day}), so that a reader such as xarray
decodes the values to dates. The record written after global step $g$ holds
\begin{equation}\label{eq:out-time}
  \tau_g=\tau_{\text{run}}+(g+1)\,\dt\,\tdir,
\end{equation}
the time at the \emph{end} of the step, in both time modes, where $\tau_{\text{run}}$ is the time at which
global step zero started. Let $g_{\text{s}}$ be the global index of the first step of the current call (zero for
a run that starts afresh). In the interpolated mode $\tau_{\text{run}}=t_{\text{first}}-g_{\text{s}}\,\dt\,\tdir$, with
$t_{\text{first}}$ the time at which the call's first step starts (the requested start date, or the time of the
start snapshot); in the discrete mode $\tau_{\text{run}}$ is the time of the snapshot $g_{\text{s}}$ records before
the call's start snapshot, the one the run started from. Under time-averaged forcing these
are the shifted midpoints (\cref{sec:input-timeavg}), so the dates are window centres, while the epoch stays
the first unshifted record. Because $\tau_g$ depends only on the global step and the start of the whole run,
the segments of a resumed chain carry exactly the labels of an uninterrupted run. In the discrete mode the
label advances by $\dt$ per step; the time-step check of \cref{sec:temp-modes} admits a small difference
between $\dt$ and the record spacing, by which the label may drift from the stored stamps. When the forcing is
cycled the time does not wrap: past the first period the dates lie beyond the end of the forcing, and an
attribute records the period. A forcing without a calendar epoch --- a bare unit, or an index axis with a
declared spacing --- gives plain seconds from its first record, with an attribute saying so.

\paragraph{Release positions and times.} For every particle the output holds its release longitude, latitude
and depth --- the seed position, after the fold into a periodic window when one is set --- and its release
time,
\begin{equation}\label{eq:out-release}
  \tau^{\text{rel}}=\tau_{\text{run}}+\max(k_{\text{rel}},0)\,\dt\,\tdir,
\end{equation}
the start of the step in which the particle is activated, with $k_{\text{rel}}$ the rounded release step of
\cref{eq:life-release} (zero without a delayed release). A particle never activated within the run still
gets its scheduled time, which then lies after the last record. The release positions travel with the
particle state and are stored in every checkpoint (\cref{sec:out-restart}), so each segment of a resumed chain
writes the same values; they are missing only for a run resumed from a checkpoint of an earlier engine.

\paragraph{Velocities.} The velocity of a record is the one the integrator evaluated for
the step that ends at the record, at the step's start position and time (for RK45, at its end); it therefore
names no coordinates. \code{vel\_w} is the rate of depth change, positive down, and on a sigma mesh the
dia-sigma flux (\cref{ch:vertical}). All components are flow velocities, with the same sign in forward and
backward runs (the integrators carry $\tdir$ times the velocity; the record removes the factor). A velocity is
written only where the particle was alive through the whole step --- alive after the release activation and
before the step, and alive after the step and the beaching hook. The record in which a particle dies or
strands, every later record of a beached or outside particle, a particle not yet released and a particle
resuspended during the step hold missing values. A two-dimensional run writes no \code{vel\_w}.

\begin{figure}[t]
\centering
\begin{tikzpicture}[scale=0.85]
  \draw[fill=gray!8]  (0.45,0.45) rectangle (4.45,2.85);
  \draw[fill=gray!12] (0.25,0.25) rectangle (4.25,2.65);
  \draw[fill=white] (0,0) rectangle (4,2.4);
  \draw[step=0.4,gray!35,line width=0.4pt] (0,0) grid (4,2.4);
  \draw[kelp,line width=1pt] (0,1.2) rectangle (1.6,2.4);
  \node[font=\footnotesize,kelp!70!black,anchor=south west] at (0,2.4) {};
  \node[font=\footnotesize] at (2,2.95) {time (records) $\longrightarrow$};
  \node[font=\footnotesize,rotate=90] at (-0.45,1.2) {trajectory $\longrightarrow$};
  \node[font=\footnotesize,anchor=west,align=left] at (4.9,1.5)
    {layers (per particle, per record):\\[2pt]
     \code{lon},\,\code{lat},\,\code{z},\,\code{status},\\
     \code{vel\_*},\,\code{failure\_reason},\,\code{temp},\,\dots\\[2pt]
     \textcolor{kelp!70!black}{one chunk: $n_t$ particles $\times$ $n_o$ records}};
\end{tikzpicture}
\caption{The output data model: two-dimensional arrays indexed by particle (trajectory) and by record (time),
one such array per stored variable, with the time and the trajectory identifier as one-dimensional coordinate
variables and the release position and time as one-dimensional arrays over the particles. The position and
status are always present, and so are the velocities of a tracking run (NaN where none was evaluated);
failure reasons and sampled scalars are added when available. In a
Zarr store each array is split into chunks of \cref{eq:out-chunks}.}
\label{fig:output}
\end{figure}

\paragraph{Provenance.} Every output carries global attributes that record how it was produced: the
integrator, the step, the mesh type, the direction, the start index and step count, whether the run is
three-dimensional and time-interpolated, the temporal reconstruction and time averaging, the engine version,
the exact-anchor flag and the resolved grid-hash cell with the rule that chose it (\cref{sec:loc-cell}); in
three dimensions the vertical-velocity combine and sign convention. Options that are off by default add their
parameters only when they are on --- pulses, diffusion, beaching, the adaptive scheme's tolerance, periodic
windows, sampling, a non-default horizontal scheme or element search, single-precision fields, read-ahead and
the rotation angles --- so that an output written with the defaults carries none of them. A run is repeated
from these attributes; the recorded grid-hash cell, for instance, can be pinned to reproduce the anchors of
an earlier run.

\section{Checkpoints and restart}\label{sec:out-restart}

When \conf{checkpoint\_freq} is set, the engine writes a checkpoint after every step $g$ for which $g+1$ is a
multiple of it, overwriting the previous one, and a final checkpoint at the end, stamped with the last step
actually executed. The checkpoint lives in the output file itself as arrays named \code{restart\_*}: the
positions, statuses and nearest nodes, the release positions, the step and forcing indices, and --- when the
run uses them --- the failure reasons, release steps, carried elements, the counts of re-wetted steps and the
frame angles. Positions and release positions are stored in the computational frame. The arrays are created
by the first checkpoint, so an output written without checkpoints holds none of them. A flag,
\code{restart\_valid}, is created before any other checkpoint array and written as zero, cleared again before
every later checkpoint, and set only after all arrays are written; a checkpoint interrupted mid-write, the
first one included, is therefore detected rather than trusted. Checkpoints are off by default: a run killed
at a wall limit without them leaves nothing to resume from.

Resuming is done by the calling program rather than by a configuration key. It loads the checkpoint
(\fn{load\_restart}), checks its frame against the mesh, and continues with the next step index, the saved
counts of re-wetted steps and the forcing position; in the interpolated mode it also supplies the simulation
time, which the checkpoint does not store. For the record times of the continuation to equal those of the
uninterrupted run (\cref{eq:out-time}), the program passes the same forcing times as the first call ---
under time-averaged forcing the shifted midpoints --- and, when a cycled forcing is read in the discrete
mode, the start index counted on without wrapping (in the discrete mode the forcing index stored in the
checkpoint is that unwrapped index, so the index after it continues the run).
The release positions need no step of their own: they come back with the particle state. A program that
splits a resumed state into parts slices it with \fn{ParticleState.subset}, which slices every per-particle
array, the release positions included. With the exact-anchor guard the nearest nodes of the resumed state
are recomputed once.

\paragraph{Resumable chains.} A run longer than one allocation of a batch system is split into a chain of
jobs. Each job resumes from the newest standalone restart file --- written with \fn{write\_restart\_file}, which
writes the same arrays to a separate NetCDF file under a temporary name and moves it into place only after a
clean close, so a failed write keeps the previous restart --- writes its records into its own output segment,
and drops a numbered restart at every checkpoint (\cref{fig:out-chain}). At the end the segments are
concatenated along time. The concatenated chain equals one uninterrupted run bit for bit, because everything
that depends on time --- the random draws of diffusion and beaching, the release activation, the record and
checkpoint cadence --- is keyed to the global step index, never to the step within a job. The library supplies
these pieces; the driver that submits the chain belongs to the experiment.

\begin{figure}[t]
\centering
\begin{tikzpicture}[x=1.15cm,y=1cm]
  \draw[-{Stealth[length=5pt]},gray] (0,0)--(10.3,0) node[right,font=\footnotesize,black]{global step $g$};
  \foreach \j/\a/\b in {1/0/3.4,2/3.4/6.8,3/6.8/10.0}{
    \draw[flowbox,fill=lagblue!8] (\a+0.06,0.35) rectangle (\b-0.06,1.05);
    \node[font=\footnotesize] at ({(\a+\b)/2},0.7) {job \j: segment \j};
  }
  \foreach \x in {0.5,1.0,...,10.0} \draw[gray!70] (\x,-0.08)--(\x,0.08);
  \foreach \x/\lbl in {3.4/restart 1,6.8/restart 2}{
    \draw[kelp,line width=1pt] (\x,-0.25)--(\x,1.3);
    \node[font=\footnotesize,kelp!70!black,anchor=south] at (\x,1.3) {\lbl};
  }
  \node[font=\footnotesize,anchor=north,align=center] at (5,-0.3) {records at $(g+1)\bmod \text{\conf{output\_freq}}=0$ continue across the job boundaries};
  \draw[flowarrow] (10.0,0.7)--(10.6,0.7) node[right,font=\footnotesize,align=left]{concatenate\\along time};
\end{tikzpicture}
\caption{A resumable chain. Each job resumes from the newest numbered restart, writes its own segment and
drops restarts as it goes; the segments are concatenated at the end. Because every cadence and every random
draw is keyed to the global step, the result equals one uninterrupted run.}
\label{fig:out-chain}
\end{figure}

\section{Merging a sharded release}\label{sec:out-merge}

A large release split into contiguous blocks of seeds and run as parallel processes (\cref{sec:perf-shard})
produces one Zarr output per shard. The merge tool concatenates them along the trajectory axis in shard
order, after checking that the shards share their time axis, their variables and every attribute that
defines the run. A run with a single release merges into an output identical in its data to the unsharded run, stochastic
diffusion and beaching included, because each shard keys its random draws by the particle's index in the
full release; the merge refuses shards passed out of that order whenever the run draws random numbers (and
warns for a deterministic run). A pulsed release is refused. Shards that are independent by design --- a
pulsed run split through its restart, one seed per shard --- merge on an explicit option that lifts the offset
order, the pulsed refusal and the seed equality, refuses shards that share any seed of a random feature they
use over the same particle indices, and records each shard's particle count, offset and seeds in the merged store. Checkpoint
arrays are not merged. The merge keeps the time attributes and the release variables, and refuses shards that
differ in the output format, in the name of the record dimension, or in the units or calendar of the time;
shards of the earlier format merge as before.

\section{Outputs of earlier engines}\label{sec:out-format1}

Outputs written before release 1.0 are in output format~1. Such a file has no
\code{kelp\_output\_format} attribute, and its record dimension is called \code{obs}. In format~1 the record time has
no epoch: in the discrete mode it is the index of the snapshot the step started from, one step behind the
position it labels, and in the interpolated mode the seconds from the first record. Format~1 has no release
variables; a two-dimensional run writes a vertical velocity of zeros; a particle's velocity is written in the
record in which it dies or strands and is zero afterwards; the discrete mode samples scalar fields at the
start of the step; and the checkpoint arrays exist in every file, holding fill values when no checkpoint was
written. These files remain readable by \fn{load\_trajectories}, \fn{load\_restart}, the merge tool and the
visualization helpers; there is no converter.

% (plan Task 5 interfaces) on the final code.
\codeanchor{The writer \fn{TrajectoryWriter} and its NetCDF and Zarr backends, the embedded checkpoint,
\fn{write\_restart\_file}, \fn{load\_restart} and \fn{load\_trajectories} are in \srcfile{writers.py}; the record
cadence and checkpoint scheduling are in the time loop of \srcfile{tracking.py}, the record time and the release
variables in \fn{\_record\_clock\_origin} and \fn{\_release\_arrays} there; the epoch and calendar of the time
come from \fn{time\_axis\_epoch} (\srcfile{readers/\_\_init\_\_.py}); the velocity mask is
\fn{mask\_inactive\_velocity} (\srcfile{core/release.py}); the frame check on resume is
\fn{\_validate\_restart\_frame} (\srcfile{tracking.py}); sharding and merging are
\code{utils/sort\_seeds.py}, \code{utils/split\_seeds.py} and \code{utils/concat\_trajectories.py}.}
